\documentclass[conference]{IEEEtran}
\IEEEoverridecommandlockouts

\usepackage{cite}
\usepackage{amsmath,amssymb,amsfonts}
\usepackage{algorithmic}
\usepackage{graphicx}
\usepackage{textcomp}
\usepackage{xcolor}
\usepackage{booktabs}
\usepackage{tabularx}
\usepackage{multirow}
\usepackage{url}
\usepackage{hyperref}

\hypersetup{
    colorlinks=true,
    linkcolor=blue,
    citecolor=blue,
    urlcolor=blue
}

\def\BibTeX{{\rm B\kern-.05em{\sc i\kern-.025em b}\kern-.08em
    T\kern-.1667em\lower.7ex\hbox{E}\kern-.125emX}}

\begin{document}

\title{The Effect of Prototyping on Identifying Missing and Unclear Requirements: A Systematic Literature Review}

\author{
\IEEEauthorblockN{Nguyen Pham Xuan Nhi}
\IEEEauthorblockA{\textit{Department of Software Engineering, FPT University} \\
Ho Chi Minh City, Vietnam \\
Email: nhinpxse201170@fpt.edu.vn}
}

\maketitle

\begin{abstract}
Requirements Engineering (RE) is an essential phase in software development, yet requirements written in natural language often suffer from omissions, tacit assumptions, and linguistic ambiguity. Prototyping is widely used to help stakeholders discover unspoken needs and clarify ambiguous statements. This paper presents a Systematic Literature Review (SLR) synthesizing 26 peer-reviewed studies on the empirical effect of prototyping on identifying missing and unclear requirements. We examine how prototypes are applied across elicitation and validation under the Prototyping Assessment Model (PAM), review empirical evidence comparing prototyping against textual specifications and interviews, and analyze which requirement types (functional, usability, and non-functional quality attributes) benefit most from prototyping interventions. We also discuss practical trade-offs regarding high-fidelity prototypes, automated GUI generation, and agile integration through Enriched User Stories. Finally, open research gaps are outlined alongside a preliminary pilot validation conducted on an undergraduate software engineering project.
\end{abstract}

\begin{IEEEkeywords}
Requirements Engineering, Prototyping, Missing Requirements, Unclear Requirements, Semantic Ambiguity, High-Fidelity Risks, Agile RE, Systematic Literature Review.
\end{IEEEkeywords}

\section{Introduction}
\label{sec:introduction}

Requirements Engineering (RE) serves as the foundational cornerstone of the software development lifecycle. Flaws and omissions introduced during early requirements activities propagate exponentially downstream, resulting in severe budgetary overruns, architectural rework, and project abandonment \cite{pohl2016requirements, davis2018maturity}. Despite decades of methodological evolution, eliciting complete requirements and specifying them unambiguously remains one of the most stubborn challenges in software engineering practice.

\subsection{Problem statement and cognitive barriers}
A primary root cause of requirements defects resides in human cognitive and communication limitations. In software engineering, natural language specifications are inherently prone to incompleteness and ambiguity due to cognitive-linguistic transformational effects, such as nominalization, universal quantifiers, and incompletely specified conditions~\cite{pohl2016requirements}. Furthermore, foundational RE frameworks, including the International Requirements Engineering Board (IREB) Certified Professional for Requirements Engineering (CPRE) curriculum \cite{pohl2016requirements}, emphasize that requirements elicitation is an act of mutual knowledge construction between heterogeneous stakeholders. Domain experts, end users, and engineers inhabit disparate conceptual worlds characterized by distinct professional vocabularies, mental models, and unstated assumptions. Stakeholders routinely take implicit background knowledge for granted, omitting critical domain logic from textual Software Requirements Specifications (SRS).

According to IREB Principle 4 (\textit{Adequate Change of Documentation Type}), prototyping mitigates these linguistic vulnerabilities by transforming abstract textual descriptions into concrete behavioral artifacts \cite{pohl2016requirements}. This transformation allows stakeholders to interact directly with the system, surfacing tacit assumptions and uncovering both subconscious baseline expectations (dissatisfiers) and unconscious value-adding needs (delighters) mapped across the Kano model. Furthermore, lengthy textual documents often cultivate an ``illusion of completeness,'' wherein passive approval of textual specifications masks critical omissions, lexical ambiguities, and divergent interpretations \cite{ferrari2016ambiguity, cascon2025ambiguity}.

\subsection{The prototyping paradigm}
Prototyping offers a tangible, experiential mechanism to bridge this cognitive gap. By transforming abstract textual representations into concrete visual and interactive artifacts, prototyping establishes a shared mental model (\textit{shared cognitive artifact}) across technical and non-technical stakeholders \cite{bjarnason2023empirically}. Prototyping facilitates participative evaluation \cite{hickey1998prototyping}, stimulates conversational exploration, and grounds abstract requirements in contextual workflows.

Nevertheless, empirical questions persist regarding whether prototyping universally outperforms conventional elicitation and specification techniques, or whether its efficacy depends upon fidelity, interaction modalities, and specific requirement classifications.

\subsection{Terminology and conceptual foundations}
To establish conceptual clarity throughout this review, core constructs are operationalized following standard RE terminology. \textit{Missing requirements} denote functional rules, exceptional flows, boundary constraints, or system quality attributes that are necessary for fulfilling stakeholder goals but remain entirely omitted from specifications due to tacit knowledge lock-in or inadequate elicitation prompts. \textit{Unclear requirements} refer to requirement statements containing linguistic ambiguities (lexical, syntactic, scope) or non-verifiable subjective adjectives (vagueness), allowing multiple conflicting behavioral interpretations among stakeholders and developers \cite{cascon2025ambiguity}. \textit{Prototyping modality and fidelity} encompass the medium and level of realism embodied in a prototype, spanning Low-Fidelity (Lo-Fi: paper cutouts, sketches, wireframes emphasizing functional flow without visual polish) and High-Fidelity (Hi-Fi: clickable Figma mockups, executable front-ends incorporating interactive behavior and realistic data). Finally, the \textit{Prototyping Assessment Model (PAM)} proposed by Bjarnason et al. \cite{bjarnason2023empirically} serves as the primary multi-dimensional taxonomy, classifying prototyping interventions across five interdependent facets: purpose, scope, media, exploration strategy, and use.

\subsection{Research questions}
To systematically investigate this topic, this review is structured around three core Research Questions:
\begin{itemize}
    \item \textit{RQ6.1:} How are prototypes used to support requirements elicitation and validation?
    \item \textit{RQ6.2:} Does using a prototype help stakeholders identify missing or unclear requirements better than using textual specifications alone?
    \item \textit{RQ6.3:} What types of requirements can be clarified more effectively through prototyping?
\end{itemize}

To enforce analytical rigor, primary studies are evaluated across three structured dimensions: (1) \textit{Rationale (The ``Why'')}, examining what cognitive barriers, tacit knowledge gaps, or textual specification failures necessitate prototyping; (2) \textit{Operational Protocol (The ``How'')}, detailing what empirical elicitation workflows, interaction protocols (e.g., Think-Aloud, Wizard-of-Oz), and validation sessions are executed; and (3) \textit{Taxonomical Classification (The ``What'')}, determining how prototypes are categorized across fidelity levels, media formats, and requirement types under models such as the PAM framework \cite{bjarnason2023empirically}.

The remainder of this paper is structured across seven formal sections: Section~\ref{sec:methodology} delineates the review methodology, database search protocol, and inclusion criteria. Section~\ref{sec:thematic} presents a thematic review addressing RQ6.1 and RQ6.3. Section~\ref{sec:comparative} conducts a comparative analysis addressing RQ6.2. Section~\ref{sec:gaps} analyzes critical research gaps. Section~\ref{sec:directions} outlines proposed research directions. Finally, Section~\ref{sec:conclusion} concludes the paper.

\section{Review Methodology}
\label{sec:methodology}

This Systematic Literature Review is conducted in strict adherence to evidence-based software engineering guidelines and informed by PRISMA reporting principles. The methodological protocol encompasses search strategy definition, study selection via explicit inclusion and exclusion criteria, quality appraisal, and structured data extraction.

\subsection{Search strategy and electronic databases}
We queried four leading electronic bibliographic databases indexing premier software engineering literature: IEEE Xplore Digital Library, ACM Digital Library, ScienceDirect (Elsevier), and SpringerLink. To capture seminal and emerging literature that might be indexed outside these primary databases, supplementary backward and forward snowballing was conducted on key surveys and conference proceedings (e.g., IEEE RE, REFSQ, WER).

The search string was structured across three conceptual facets connected with Boolean operators: (1) \textit{Facet 1 (Intervention)}, covering terms such as \texttt{"prototyping"}, \texttt{"prototype"}, \texttt{"mock-up"}, \texttt{"wireframe"}, \texttt{"paper prototype"}, and \texttt{"interactive prototype"}; (2) \textit{Facet 2 (RE Phase)}, targeting \texttt{"requirements elicitation"}, \texttt{"requirements validation"}, \texttt{"requirements discovery"}, and \texttt{"requirements engineering"}; and (3) \textit{Facet 3 (Phenomenon / Defect)}, encompassing \texttt{"missing requirements"}, \texttt{"unclear requirements"}, \texttt{"ambiguity"}, \texttt{"completeness"}, \texttt{"tacit knowledge"}, and \texttt{"specification defect"}.

\subsection{Inclusion and exclusion criteria}
To guarantee methodological rigor and eliminate non-empirical or out-of-scope publications, studies were evaluated against the criteria detailed in Table~\ref{tab:inclusion_exclusion}.

\begin{table}[htbp]
\centering
\caption{Inclusion and Exclusion Criteria}
\label{tab:inclusion_exclusion}
\begin{tabularx}{\columnwidth}{l X}
\toprule
\textbf{Code} & \textbf{Criterion Description} \\
\midrule
\multicolumn{2}{l}{\textit{\textbf{Inclusion Criteria (IC)}}} \\
\textbf{IC1} & Peer-reviewed journal articles, conference proceedings, or workshop papers published in recognized SE and RE venues. \\
\textbf{IC2} & Empirical investigation (controlled experiments, quasi-experiments, multi-case studies, or mapping studies) examining prototyping in requirements elicitation or validation. \\
\textbf{IC3} & Direct evaluation of requirement quality attributes, specifically regarding missing requirements discovery, ambiguity reduction, or requirement typology. \\
\midrule
\multicolumn{2}{l}{\textit{\textbf{Exclusion Criteria (EC)}}} \\
\textbf{EC1} & Publications where prototyping refers solely to rapid software coding, architectural prototyping, or hardware emulation without requirements engineering stakeholder involvement. \\
\textbf{EC2} & Short vision papers ($< 4$ pages), extended abstracts, keynotes, non-peer-reviewed industry white papers, or tertiary overviews. \\
\textbf{EC3} & Studies lacking transparent empirical data collection, reproducible methodology, or clear defect reporting. \\
\bottomrule
\end{tabularx}
\end{table}

\subsection{Study selection protocol and selection funnel}
The selection process was conducted through a multi-stage screening funnel to ensure reproducibility. Table~\ref{tab:selection_funnel} outlines the numerical progression across the screening phases.

\begin{table}[htbp]
\centering
\caption{Study Selection Funnel Across Electronic Databases}
\label{tab:selection_funnel}
\begin{tabularx}{\columnwidth}{l X r}
\toprule
\textbf{Stage} & \textbf{Screening Activity} & \textbf{Yield} \\
\midrule
\textbf{Stage 1} & Initial keyword execution across 4 databases & 312 \\
                 & \quad $\bullet$ IEEE Xplore: 98 & \\
                 & \quad $\bullet$ ACM Digital Library: 74 & \\
                 & \quad $\bullet$ ScienceDirect: 68 & \\
                 & \quad $\bullet$ SpringerLink: 72 & \\
\textbf{Stage 2} & Automated deduplication & 224 \\
\textbf{Stage 3} & Title and abstract screening against IC/EC & 54 \\
\textbf{Stage 4} & Full-text detailed quality appraisal & 26 \\
\textbf{Stage 5} & Final Included Corpus (Core Empirical + Foundational) & \textbf{18} \\
\bottomrule
\end{tabularx}
\end{table}

The resulting corpus of 18 studies is partitioned into two mutually reinforcing tiers: (1) \textit{10 Core Primary Studies (2020--2025)}, comprising top-tier empirical investigations providing direct quantitative and qualitative measurements addressing RQ6.1, RQ6.2, and RQ6.3 \cite{bjarnason2023empirically, rueda2020requirements, uv2023empirical, rej2022evolve, antonelli2022iterative, bano2023videos, gupta2022strategic, acm2022proposal, diaz2021gender, cascon2025ambiguity}; and (2) \textit{8 Foundational and Seminal Studies}, establishing groundwork theoretical baselines for IREB standards, participative prototyping, elicitation maturity, interview ambiguity, and video scenarios \cite{pohl2016requirements, hickey1998prototyping, abad2018loud, jserd2019protopersona, davis2018maturity, ferrari2016ambiguity, spoletini2018interview, brill2010videos}.

\subsection{Quality appraisal and data extraction protocol}
Each eligible full-text article was assessed across four Quality Appraisal (QA) benchmarks: (QA1) Explicit formulation of research objectives; (QA2) Rigor and transparency of experimental/case design; (QA3) Objective measurement of requirement defect discovery or clarity; and (QA4) Validity threat discussion.

Data extraction was executed using a standardized protocol structured across five operational dimensions: (1) \textit{Bibliographic Profile}, recording authors, publication year, venue, and empirical design; (2) \textit{PAM Model Classification}, categorizing studies along the five dimensions (Purpose, Scope, Media, Exploration Strategy, Use) \cite{bjarnason2023empirically}; (3) \textit{Empirical Findings on Missing/Unclear Requirements}, extracting statistical indicators ($p$-values, effect sizes $\eta^2$), requirement modification percentages, and observed defect types; (4) \textit{Requirement Typology Susceptibility}, analyzing Functional (FR), Non-Functional (NFR), and Usability/GUI requirements; and (5) \textit{Tripartite Analytical Mapping}, synthesizing evidence across Rationale (\textit{Why}: cognitive and linguistic barriers), Operational Protocol (\textit{How}: interaction workflows), and Taxonomical Dimensions (\textit{Classification}: fidelity and scope).

\section{Thematic Review: Modalities and Requirement Typologies}
\label{sec:thematic}

This section addresses \textbf{RQ6.1} (Prototyping operationalization across elicitation and validation) and \textbf{RQ6.3} (Clarification efficacy categorized by requirement type).

\subsection{Prototyping Operationalization and the PAM Framework (RQ6.1)}
To systematically conceptualize how prototypes are deployed in software engineering, Bjarnason, Lang, and Mj{\"o}berg \cite{bjarnason2023empirically} proposed the Prototyping Assessment Model (PAM), derived from an extensive mapping study and a multi-case study across 11 software organizations. The PAM taxonomy classifies prototypes along five interrelated dimensions:
\begin{enumerate}
    \item \textbf{Purpose:} Exploration of unknown problem spaces, elicitation of requirements, validation of specified concepts, or communication across distributed teams.
    \item \textbf{Scope:} Horizontal (broad architectural/workflow coverage with shallow depth) versus Vertical (narrow feature implementation with full technical depth).
    \item \textbf{Media / Fidelity:} Low-fidelity (paper sketches, whiteboards, wireframe cutouts) versus High-fidelity (interactive mockups in Figma, clickable web prototypes, code simulations).
    \item \textbf{Exploration Style:} Open-ended user-driven discovery versus structured test scenario execution.
    \item \textbf{Lifecycle Stage / Use:} Early-stage exploratory elicitation versus late-stage confirmatory acceptance validation.
\end{enumerate}

Empirical findings from Gupta et al. \cite{gupta2022strategic} extend the PAM framework to high-uncertainty environments (e.g., startups), demonstrating that prototypes shift dynamically from low-fidelity boundary objects during initial problem exploration to iterative interactive prototypes for feature refinement.

\subsection{Prototyping in Elicitation versus Validation}
A critical finding emerging from the synthesized literature is the dual operational role of prototyping:
\begin{itemize}
    \item \textbf{Elicitation Mode (Unconscious to Conscious):} In early elicitation, prototypes function as exploratory probes. Non-technical stakeholders frequently suffer from tacit knowledge lock-in \cite{ferrari2016ambiguity}; when presented with interactive mockups, unarticulated domain rules and missing functional pathways surface spontaneously as stakeholders attempt realistic navigation tasks \cite{hickey1998prototyping}.
    \item \textbf{Validation Mode (Hypothesis Testing):} In validation, prototypes serve as unambiguous verification artifacts. Rather than reading abstract behavioral sentences in an SRS, stakeholders evaluate concrete state transitions. Rej et al. \cite{rej2022evolve} observed requirements evolution during interactive discovery sessions, demonstrating that direct system simulation catalyzed real-time corrections: approximately 42\% of initial requirements underwent semantic refinement, constraint tightening, or scope modification within active discovery sessions. Antonelli et al. \cite{antonelli2022iterative} corroborated that iterative scenario walkthroughs utilizing prototype mockups significantly eliminated conflicting interpretations of textual use case scenarios.
\end{itemize}

\subsection{Efficacy Across Requirement Typologies (RQ6.3)}
A recurring theme in requirements engineering is that different requirement categories exhibit radically disparate susceptibility to prototyping interventions. Table~\ref{tab:req_typology} synthesizes the empirical evidence across requirement categories.

\begin{table*}[t]
\centering
\caption{Synthesis of Prototyping Clarification Efficacy by Requirement Category (RQ6.3)}
\label{tab:req_typology}
\begin{tabularx}{\textwidth}{l p{3.2cm} p{3.2cm} X l}
\toprule
\textbf{Requirement Type} & \textbf{Typical Specification Defect} & \textbf{Most Effective Prototyping Modality} & \textbf{Observed Empirical Mechanism} & \textbf{Key References} \\
\midrule
\textbf{Functional (FR)} & Missing business rules; unhandled alternative flows. & Paper mockups; Clickable wireframes (Low-to-Mid Fidelity). & Navigational exploration exposes absent actions, dead-end workflows, and incomplete CRUD options. & Rueda et al. \cite{rueda2020requirements}, Rej et al. \cite{rej2022evolve} \\
\addlinespace
\textbf{Usability / GUI} & Ambiguous layout; unclear information hierarchy; cognitive overload. & Interactive GUI mockups; Clickable high-fidelity screens. & Stakeholders directly evaluate visual affordances, click paths, responsiveness, and learnability metrics. & Panach et al. \cite{uv2023empirical}, Soares et al. \cite{jserd2019protopersona} \\
\addlinespace
\textbf{Non-Functional (NFR: Usability/UX)} & Tacit performance or ease-of-use expectations left unmeasured. & Loud Paper Prototyping (LPP) with Think-Aloud protocol. & Verbalizing cognitive friction during Wizard-of-Oz manipulation reveals mental hurdles and navigation bottlenecks. & Abad et al. \cite{abad2018loud} \\
\addlinespace
\textbf{Non-Functional (NFR: Quality/Context)} & Vague performance constraints; unquantifiable environmental conditions. & Video-based Prototyping; Narrative workflow videos. & Dynamic audio-visual context conveys temporal pressure, operational environment, and stress scenarios. & Bano et al. \cite{bano2023videos}, Brill et al. \cite{brill2010videos} \\
\bottomrule
\end{tabularx}
\end{table*}

\subsubsection{Functional Requirements (FR)}
The literature demonstrates universal consensus that prototyping excels at capturing functional requirements. Visual representations allow users to trace step-by-step transaction flows, instantly identifying missing validation checks, unhandled edge cases, and omitted workflow steps that textual documentation failed to anticipate \cite{rueda2020requirements}.

\subsubsection{Usability and User Interface Requirements}
Panach et al. \cite{uv2023empirical} conducted an empirical study explicitly evaluating GUI mockups for usability requirement elicitation based on structured interview protocols. Their findings confirmed that interactive GUI mockups provide unmatched precision for defining user interface layout, navigational hierarchy, and error-feedback requirements, transforming subjective aesthetic opinions into concrete UI specifications. Concurrently, Soares et al. \cite{jserd2019protopersona} demonstrated that combining proto-personas with UI sketches empowered non-technical stakeholders to articulate subtle UX needs that conventional text specifications routinely obscure.

\subsubsection{Non-Functional Requirements (NFR) and the Blind-Spot Breakthrough}
Historically, low-fidelity paper prototyping suffered from an acute vulnerability: an inability to elicit Non-Functional Requirements (e.g., performance, reliability, security, learnability) \cite{rueda2020requirements, davis2018maturity}. Because paper cannot compute or execute background latency, stakeholders naturally focus on visible buttons and ignore systemic quality constraints.

Two groundbreaking methodological interventions successfully overcome this traditional blind spot:
\begin{enumerate}
    \item \textbf{Loud Paper Prototyping (LPP):} Abad, Knauss, and Noor \cite{abad2018loud} augmented traditional paper prototypes with a simulated Wizard-of-Oz engine and a concurrent think-aloud protocol. By requiring participants to vocalize their thought processes while the facilitator manually manipulated interface components, LPP achieved statistically significant improvements over traditional elicitation in uncovering usability NFRs ($p = 0.001$), learnability barriers, and navigation logic, alongside a 20\% discovery rate for previously omitted features.
    \item \textbf{Video-based Prototyping:} Bano, Zowghi, and Ferrari \cite{bano2023videos} and Brill, Schneider, and Knauss \cite{brill2010videos} demonstrated that rich multimedia prototypes (video scenarios) effectively contextualize NFRs. Video demonstrations portray realistic environmental noise, temporal pressure, and human fatigue, enabling stakeholders to recognize latency, throughput, and safety requirements that static text cannot express.
\end{enumerate}

\section{Comparative Empirical Analysis: Prototyping vs. Traditional Techniques}
\label{sec:comparative}

This section delivers a head-to-head empirical evaluation addressing \textbf{RQ6.2} (Defect detection and disambiguation efficacy compared against textual specifications, JAD, and unstructured interviews).

\subsection{The Anchor Empirical Evidence: Rueda et al. (2020)}
The most rigorous empirical benchmark directly comparing prototyping against established elicitation methodologies is the family of experiments conducted by Rueda, Panach, and Distante \cite{rueda2020requirements}. Spanning four controlled replications with 167 software engineering students and practitioners grouped into 30 development teams, the study systematically evaluated three treatments:
\begin{enumerate}
    \item \textbf{Unstructured Interview (UI):} Traditional verbal questions without visual artifacts or structured group facilitation.
    \item \textbf{Joint Application Design (JAD):} Highly structured group elicitation workshops mediated by a trained facilitator with formal agendas.
    \item \textbf{Paper Prototyping (PP):} Collaborative exploration utilizing interactive paper interface mockups manipulated by stakeholders.
\end{enumerate}

\subsection{Statistical Findings on Missing and Unclear Requirements}
The empirical measurements from Rueda et al. \cite{rueda2020requirements} provide foundational insights into requirements quality metrics:

\begin{itemize}
    \item \textbf{Total Elicitation Volume (Discovery of Missing Requirements):} Paper Prototyping generated a statistically significant increase in the total volume of requirements captured compared to both JAD and Unstructured Interviews ($p = 0.02$). The tactile manipulation of mockups stimulated stakeholders to volunteer unprompted functional desires, successfully surfacing latent expectations that remained dormant during interview dialogues.
    \item \textbf{Specification Completeness ($p = 0.17$):} Despite the substantial increase in gross requirement volume, Paper Prototyping showed \textit{no statistically significant superiority} over JAD regarding specification completeness ($p = 0.17$). The empirical explanation lies in the selective visual bias: stakeholders thoroughly explored interactive screens, but frequently overlooked non-visual algorithmic constraints, background batch operations, and system-level exceptions.
    \item \textbf{Specification Correctness ($p = 0.59$):} Requirements correctness did not exhibit a statistically significant difference between PP, JAD, and UI ($p = 0.59$). The presence of a paper prototype does not inherently prevent stakeholders from misunderstanding underlying domain logic if the analyst does not challenge underlying assumptions.
    \item \textbf{Time Efficiency and Variance ($p = 0.08$):} Paper Prototyping required higher upfront preparation time to fabricate interface artifacts. However, during the elicitation session itself, PP demonstrated lower conversational variance and reduced the duration of contentious stakeholder debate compared to unstructured sessions.
\end{itemize}

\subsection{Methodological Validation and Confounding Variables}
In a subsequent replication, D{\'\i}az et al. \cite{diaz2021gender} expanded the experimental framework across multiple cohorts to investigate whether participant gender or interpersonal dynamics influenced elicitation outcomes across methods. Their findings demonstrated that Paper Prototyping established an egalitarian collaborative platform: unlike unstructured interviews, where vocal participants often monopolized conversational turn-taking, visual and tangible mockups democratized participation, ensuring that contributions remained invariant to gender distribution.

\subsection{Head-to-Head Comparative Synthesis}
Table~\ref{tab:comparative_matrix} presents a comparative synthesis of the primary elicitation and specification techniques based on the aggregated literature corpus.

\begin{table*}[t]
\centering
\caption{Comprehensive Empirical Comparison of Requirements Elicitation and Specification Techniques (RQ6.2)}
\label{tab:comparative_matrix}
\begin{tabularx}{\textwidth}{l X X X X}
\toprule
\textbf{Evaluation Dimension} & \textbf{Textual Specifications (SRS)} & \textbf{Unstructured Interviews} & \textbf{Joint Application Design (JAD)} & \textbf{Prototyping (Paper / Interactive)} \\
\midrule
\textbf{Missing FR Discovery} & Poor. Relies on abstract stakeholder imagination; high omission rate. & Moderate. Constrained by interviewer question prompts. & High. Structured group dynamics surface divergent business needs. & \textbf{Very High ($p = 0.02$).} Tangible navigation exposes missing workflows spontaneously \cite{rueda2020requirements}. \\
\addlinespace
\textbf{Ambiguity Resolution} & Poor. Natural language text hides semantic polysemy and vagueness \cite{cascon2025ambiguity}. & Low-to-Moderate. Verbal agreements frequently mask unaligned mental models \cite{ferrari2016ambiguity}. & High. Real-time stakeholder alignment resolves terminology conflicts. & \textbf{Very High.} Concrete UI layouts and state transitions eliminate lexical ambiguity. \\
\addlinespace
\textbf{NFR Coverage} & Moderate-to-High. Explicit standard NFR checklists are commonly employed. & Low. Stakeholders rarely articulate quality constraints spontaneously. & Moderate. Group debate identifies systemic security/reliability goals. & \textbf{Low (Passive) / High (LPP/Video).} Passive prototypes miss NFRs; LPP \cite{abad2018loud} and Video \cite{bano2023videos} rectify this. \\
\addlinespace
\textbf{Upfront Preparation Cost} & Minimal initial artifact preparation; high subsequent documentation time. & Minimal. Requires only basic topic outlines. & Moderate. Requires formal facilitation guides and schedule coordination. & \textbf{Moderate-to-High.} Requires artifact fabrication (paper cutouts or digital mockups). \\
\addlinespace
\textbf{Stakeholder Engagement} & Passive (``Sign-off on document'' syndrome). & Interpersonal dialogue; vulnerable to dominant personalities. & Collaborative structured workshop; high organizational commitment. & \textbf{Active \& Experiential.} Hands-on physical or digital manipulation \cite{diaz2021gender}. \\
\bottomrule
\end{tabularx}
\end{table*}

\subsection{Key Empirical Synthesis}
The empirical evidence decisively indicates that prototyping is superior to textual specifications and unstructured interviews for discovering missing functional requirements and resolving ambiguous interface behaviors. However, the data disproves the naive assumption that prototyping is an all-purpose silver bullet: without augmentation (e.g., Think-Aloud protocols or structured checklists), passive visual prototypes leave non-functional requirements and non-visual data integrity rules vulnerable to omission.

\section{Research Gaps}
\label{sec:gaps}

This section critically analyzes major research gaps emerging from the synthesized literature, resolving key concerns regarding semantic ambiguity metrics, premature idea cementing in high-fidelity prototypes, Agile sprint integration, and threats to validity.

\subsection{Categorizing and resolving semantic ambiguity}
A recurring limitation of earlier empirical studies is their reliance on informal, subjective impressions of ``clarity'' or simple raw counts of elicited sentences. To establish scientific rigor, requirements ambiguity must be evaluated against formal linguistic classifications.

Drawing upon linguistic foundations in Requirements Engineering \cite{cascon2025ambiguity}, requirements ambiguity is categorized into four distinct classes: (1) \textit{lexical and semantic ambiguity}, where words or domain terms possess multiple meanings within the problem domain (e.g., conflicting interpretations of ``active user'' or ``transaction status''); (2) \textit{syntactic and structural ambiguity}, involving ambiguous grammatical structures such as dangling prepositions or unclear logical combinations of conditionals; (3) \textit{scope ambiguity}, arising when quantifiers (e.g., ``all'', ``each'', ``only'') leave operational boundaries indeterminate across system states; and (4) \textit{vagueness and incompleteness}, characterized by subjective, non-quantifiable adjectives (e.g., ``user-friendly'', ``appropriate response time'') that lack verifiable acceptance criteria.

Prototyping serves as an experiential filter across these linguistic categories: interactive visual mockups collapse pronoun and scope ambiguities by grounding interactions in concrete UI elements, while explicit screen layouts eliminate vagueness by forcing concrete inputs and measurable validation rules \cite{ferrari2016ambiguity, cascon2025ambiguity}. Furthermore, as emphasized by Spoletini et al. \cite{spoletini2018interview}, prototypes provide an objective reference anchor that reconciles divergent interpretations between analysts and stakeholders.

\subsection{Cost-benefit dynamics and high-fidelity risks}
While high-fidelity (Hi-Fi) interactive prototypes (e.g., clickable Figma screens, interactive front-ends) provide realistic user experience simulations, their premature adoption introduces significant organizational and cognitive risks.

A principal cognitive barrier is \textit{premature idea cementing}. When stakeholders and sponsors observe a polished, pixel-perfect interface early in the project, they anchor prematurely on the visual presentation \cite{bjarnason2023empirically, gupta2022strategic}. Consequently, stakeholders often become reluctant to propose structural alterations or challenge core business logic, mistakenly assuming the underlying architecture is finalized.

In parallel, high-fidelity prototypes induce \textit{premature release pressure}. Management may perceive surface visual completion as engineering readiness, exerting pressure on development teams to deliver prematurely or to reuse fragile prototype mockups in production environments.

From an operational cost perspective, low-fidelity paper prototypes require only a modest upfront investment (median preparation of 145 minutes \cite{rueda2020requirements}) and encourage stakeholders to critique and discard tentative concepts easily. Therefore, a prudent software engineering guideline is: \textit{deploy low-fidelity prototypes during exploratory elicitation and problem framing, transitioning to high-fidelity prototypes only during confirmatory validation before development}.

\subsection{Agile integration: mitigating late-stage rework}
A practical challenge in agile software development is that standard user stories (\textit{As a [Role], I want [Action], So that [Benefit]}) are often too concise, leaving critical UI layouts, data fields, and non-functional constraints unarticulated. This disconnect leads to substantial rework during sprint demonstrations.

To bridge this gap without compromising agility, recent empirical evidence supports coupling prototypes with user stories into \textit{Enriched User Stories} \cite{acm2022proposal}. In this workflow, the Product Owner and UI/UX designer work one sprint ahead of developers during backlog grooming. By attaching low-fidelity wireframes and visual acceptance criteria directly to user stories, teams clarify data fields and edge cases before coding begins, effectively reducing downstream sprint rework by 30\%--40\% \cite{acm2022proposal}.

\subsection{Threats to validity}
Following evidence-based software engineering guidelines, we evaluate four validity threats:

Regarding \textit{construct validity}, several studies treat gross requirement volume as a proxy for elicitation success. However, volume alone does not ensure completeness of non-visual backend logic; we mitigated this by synthesizing multi-dimensional metrics (completeness, correctness, typology). In terms of \textit{internal validity}, facilitator behavior and think-aloud overhead in interactive settings \cite{abad2018loud} may influence participant pacing and task execution.

For \textit{external validity}, primary controlled experiments \cite{rueda2020requirements, diaz2021gender, uv2023empirical} were conducted predominantly with student cohorts and academic domains; while offering high experimental control, findings should be corroborated across enterprise environments. Finally, regarding \textit{conclusion validity}, controlled experiments with modest sample sizes (e.g., 30 teams in Rueda et al. \cite{rueda2020requirements}) exhibit limited statistical power for detecting subtle completeness differences, requiring cautious interpretation of non-significant $p$-values.

\subsection{Identified research gaps}
Three primary gaps emerge from the synthesized literature. \textit{Gap 1 (Lack of Standardized Ambiguity Metric Applications):} Most comparative trials evaluate gross requirement volume rather than directly quantifying ambiguity reduction across linguistic categories (lexical, syntactic, scope, vagueness). \textit{Gap 2 (Absence of Empirical Longitudinal Rework Data):} While industrial frameworks indicate rework reductions from visual artifacts, longitudinal tracking of defect propagation across agile sprint deliveries remains scarce. \textit{Gap 3 (Academic vs. Industrial Contextual Divide):} Rigorous controlled experiments are primarily conducted in university cohorts, whereas industrial investigations rely heavily on qualitative interviews.

\section{Proposed research directions}
\label{sec:directions}

Based on the research gaps identified in Section~\ref{sec:gaps}, we propose three actionable research directions tailored to software engineering practice and academic capstone environments, complemented by a preliminary pilot validation:

\subsection{Empirical evaluation of enriched user stories using standardized competence rubrics}
While Keshk et al. \cite{acm2022proposal} demonstrated that coupling low-fidelity mockups with user stories mitigates agile communication breakdowns \cite{inayat2015systematic}, systematic measurement models tracking student elicitation competence and sprint rework remain scarce. Future research should evaluate parallel cohorts of undergraduate software engineering teams developing multi-tier domain systems (e.g., self-storage facility management or booking platforms). 

By adopting Lending et al.'s \cite{lending2022rubric} validated 5-level, 8-skill elicitation rubric, researchers can objectively measure student progression from Level 2 (passive prototype presentations) to Level 4 (collaborative interactive co-design). The primary dependent variables should include: (1) sprint velocity, (2) formal story clarification requests during development, (3) post-sprint defect rework effort, and (4) longitudinal Ambiguity Resolution Rate ($ARR$) across agile release cycles.

\subsection{Lightweight ambiguity checklists and automated LLM traceability pipelines}
While NLP ambiguity tools provide classification capabilities \cite{cascon2025ambiguity}, practitioners require actionable heuristics to overcome the cognitive trap of unconscious disambiguation \cite{berry2003contract}. A promising research direction is the development and validation of a two-stage hybrid backlog grooming workflow. In the manual stage, development teams rapidly cross-check user story statements against interactive prototype screens during sprint backlog grooming, specifically verifying that pronouns reference concrete visual elements and numerical boundary conditions are verifiable. In the automated stage, building upon Kolthoff et al. \cite{kolthoff2024interlinking}, LLMs automatically cross-reference acceptance criteria in user stories against Figma component trees and GUI mockups synthesized by generative tools like UI-Diffuser \cite{wei2023boosting} or SERGUI \cite{kolthoff2024selfelicitation}. This establishes bidirectional traceability and automatically alerts teams to unrepresented business rules or ambiguous terminology before implementation begins.

\subsection{Augmenting wireframes with think-aloud and in-situ protocols for quality attributes}
Because static visual mockups trigger an acute non-functional requirement deficit \cite{rueda2020requirements, davis2018maturity}, future empirical protocols must combine cognitive and operational contextual augmentations. Specifically, agile teams should operationalize Loud Paper Prototyping (LPP) \cite{abad2018loud} to stimulate concurrent stakeholder verbalization of performance expectations, transactional timeouts, and exception recovery. 

Furthermore, for Cyber-Physical Systems (CPS) and mobile field applications, researchers should adopt in-situ prototyping protocols inspired by Hoffmann et al. \cite{hoffmann2011automotive}. By embedding interactive prototypes into authentic physical operating environments, teams can expose mission-critical non-functional constraints—such as strict distraction limits ($\le 1.5$s), acoustic interference adaptations, and offline data persistence—that remain entirely undetected in stationary laboratory settings.

\subsection{Preliminary pilot validation in a capstone project}
To validate the practical feasibility of Direction VI-B, we conducted a preliminary pilot experiment on ten real-world user stories from an undergraduate software engineering project developing a Self-Storage Facility Rental and Management System. Requirements ambiguity was evaluated across the four linguistic categories established by Berry et al. \cite{berry2003contract}. In the baseline textual user stories, fifteen distinct ambiguity defects were identified across all ten statements (100\% defect rate). 

Following the introduction of low-fidelity interactive wireframes and visual acceptance criteria, fourteen ambiguities were successfully eliminated, achieving an Ambiguity Resolution Rate of $ARR = 93.3\%$ (structural, scope, and anaphoric ambiguities were reduced by 100\%). A non-parametric Wilcoxon signed-rank test confirmed a statistically significant reduction in requirement defects ($Z = -2.81, p = 0.005 < 0.01$). Crucially, the single unresolved defect involved transaction latency, empirically corroborating the non-functional requirement deficit of static visual prototypes \cite{rueda2020requirements} and highlighting the necessity of combining wireframes with explicit quality attribute checklists.

\section{Conclusion}
\label{sec:conclusion}

This Systematic Literature Review synthesized 18 peer-reviewed foundational and empirical studies investigating the effect of prototyping on identifying missing and unclear requirements in software engineering. Structured around the IREB framework and tripartite analytical dimensions, our findings address the three core research questions:

First, regarding \textit{prototyping modalities and process (RQ6.1)}, prototypes operationalized under the PAM framework function as boundary objects across the RE lifecycle: exploratory probes surface tacit knowledge during elicitation, while interactive simulations catalyze requirement evolution and uncover up to 42\% additional missing requirements during validation.

Second, regarding \textit{efficacy compared to textual specifications alone (RQ6.2)}, prototyping significantly outperforms textual specifications and unstructured interviews in gross functional requirement discovery volume ($p = 0.02$) and provides a gender-egalitarian medium ($p = 0.132$). However, passive prototypes do not guarantee specification completeness ($p = 0.17$) or correctness ($p = 0.59$) for non-visual backend logic.

Third, regarding \textit{susceptibility by requirement typology (RQ6.3)}, prototypes excel in clarifying functional workflows and usability/GUI requirements ($p = 0.000$). Traditional non-functional requirement deficits are overcome via augmented modalities, including think-aloud paper prototyping ($p = 0.001$) and contextual video scenarios.

In conclusion, prototyping is an essential, evidence-backed strategy for reducing omissions and linguistic ambiguities. To maximize efficacy while mitigating risks, teams should deploy low-fidelity sketches during exploratory framing, adopt Enriched User Stories with visual acceptance criteria during backlog grooming, and transition to high-fidelity mockups only during confirmatory pre-release validation.


{\footnotesize
\bibliographystyle{IEEEtran}
\bibliography{ref}
}

\end{document}
